%!TeX root=./tese.tex
% Related work. Sources verified against the local paper library
% (/home/gio/mestrado/openclaw/data/papers*), the rag corpus, and the Neo4j
% citation graph (in-library citation counts guided selection); recent
% works additionally verified by web search. \emph{[Placeholder: ...]}
% items to resolve before the final build.
\chapter{Related Work}
\label{cap:related}

Five bodies of work border this thesis: serving systems for large language
models, serving systems for \emph{multimodal} models, autoscalers for
model serving, workload characterization studies, and time-series
forecasting models. This chapter reviews each through the lens the thesis
argues from --- \emph{which load signal a system watches, and what it does
with it} --- and closes with a positioning summary. Throughout, the
provisioning vocabulary is the one fixed in
Chapter~\ref{cap:intro}: a \emph{pool} is an independently sized set of
workers for one inference stage, and \emph{scaling} means changing a
pool's worker count over time under latency service-level objectives
(SLOs).

\section{Serving large language models}
\label{sec:rel-llm}

Modern LLM serving rests on engine-level scheduling innovations. Orca
introduced iteration-level scheduling --- admitting and evicting requests
at token-iteration granularity rather than request
granularity~\cite{Yu2022Orca}; vLLM's PagedAttention virtualized the KV
cache to eliminate memory fragmentation, making large continuous batches
practical~\cite{Kwon2023vLLM}; SGLang added structured-program execution
and prefix-cache reuse~\cite{Zheng2024SGLang}; Sarathi-Serve piggybacks
chunked prefills onto decode batches to tame the throughput--latency
trade-off~\cite{Agrawal2024SarathiServe}; and FastServe schedules
preemptively against head-of-line blocking~\cite{Wu2023FastServe}. All of
these operate \emph{within} a fixed fleet: they decide which request runs
where and when, not how many workers should exist.

A second line disaggregates the two phases of LLM inference themselves.
Splitwise~\cite{Patel2024Splitwise} and DistServe~\cite{Zhong2024DistServe}
place prefill and decode on separate worker pools, sized and scheduled to
their own bottlenecks, and demonstrate substantial goodput-per-GPU gains;
Mooncake operates this split in production at Moonshot AI, organizing the
cluster around a disaggregated KV cache and adding a prediction-based
early-rejection policy for overload~\cite{Qin2025Mooncake}; DynaServe
generalizes the split into an elastic spectrum between aggregated and
disaggregated execution~\cite{Ruan2025DynaServe}. NVIDIA Dynamo packages
this architecture as a datacenter-scale framework --- KV-aware routing,
tiered KV-cache management, disaggregated worker pools, and an SLA-driven
planner~\cite{Dynamo2025} --- and is the platform this thesis builds its
simulation on. This line establishes the premise the thesis starts from:
heterogeneous inference stages deserve separately provisioned pools. Its
control planes, however, schedule and admit \emph{within} whatever fleet
exists; the question of how each pool's size should track demand over
time is delegated to an autoscaler --- and, as
Section~\ref{sec:rel-autoscale} shows, that autoscaler watches
request-level signals.

\section{Serving multimodal models}
\label{sec:rel-mm}

Multimodal inference prepends an encoding stage to prefill and decode, and
a recent line of systems disaggregates it. ModServe, the system whose
production trace this thesis analyzes, decouples image encoding from
language-model execution at cluster scale and reports 25--41\% GPU-hour
savings from modality-aware placement and scaling on Azure
traffic~\cite{Qiu2025ModServe}. EPD disaggregation moves
encode, prefill, and decode onto dedicated resources, with asynchronous
transfer of multimodal tokens between stages, reporting large gains in
memory headroom, supported batch sizes, and TTFT SLO
attainment~\cite{Singh2025EPD}. HydraInfer pursues the same
three-stage split with hybrid scheduling of encode and decode
instances~\cite{Dong2025HydraInfer}, and vLLM-Omni extends stage
disaggregation to any-to-any multimodal pipelines that mix autoregressive
and diffusion components~\cite{vLLMOmni2026}. ElasticMM makes the encoder
pool's elasticity explicit, organizing MLLM serving into modality-decoupled
resource groups that are scaled at runtime~\cite{Liu2025ElasticMM}; RServe
occupies the opposite design point, overlapping encoding with prefill
inside a single instance rather than separating them~\cite{Guo2025RServe};
and cross-tier deployments that place the compute-bound encode stage on
consumer-class GPUs --- the hardware split this thesis's encoder pool
adopts --- have been costed explicitly~\cite{Yu2026HeteroServe}. Encoder
disaggregation has
also reached the mainstream engines: both vLLM and SGLang now ship
encoder-separated serving modes with elastic encoder
scaling.\footnote{vLLM, \emph{Encoder Disaggregation for Scalable
Multimodal Model Serving},
\texttt{https://vllm.ai/blog/2025-12-15-vllm-epd}; LMSYS, \emph{EPD
Disaggregation: Elastic Encoder Scaling for Vision-Language Models in
SGLang}, \texttt{https://www.lmsys.org/blog/2026-01-12-epd/}. Both
verified 2026-07-07.}

These systems settle the architecture question --- the encoder pool
exists, separately deployed, separately scalable, and now shipping in
mainstream engines --- and thereby make the question this thesis answers
operational rather than hypothetical. Their scaling decisions are
driven by currently observed load (queue depths, utilization, arrival
rates), inheriting the reactive loop of Section~\ref{sec:rel-autoscale};
none of them asks \emph{what demand signal} an encoder pool should be
provisioned on, whether that signal is visible in the request stream a
conventional autoscaler watches (Chapter~\ref{cap:trace-eda} shows it is
not), or whether it is forecastable at the horizons replica cold starts
require (Chapter~\ref{cap:predictors} shows it is). ModServe's own trace
analysis documents bursty, heterogeneous image traffic --- the motivating
observation --- but its scaling remains reactive; the per-pool demand
signals of its released trace were, to the best of our knowledge, not
characterized as forecasting targets before this thesis.

\section{Autoscaling for model serving}
\label{sec:rel-autoscale}

SLO-aware autoscaling predates LLMs. Swayam derived provisioning rules
from SLA models for ML inference services~\cite{Gujarati2017Swayam}; MArk
combined proactive and reactive scaling over heterogeneous cloud
instances~\cite{Zhang2019MArk}; Clockwork built latency predictability
bottom-up, making worker performance a modelable
quantity~\cite{Gujarati2020Clockwork}; INFaaS automated model-variant
selection and scaling behind a model-less
interface~\cite{Romero2021INFaaS}; AlpaServe showed that statistical
multiplexing across model-parallel replicas changes the provisioning
calculus~\cite{Li2023AlpaServe}. Closest in spirit to per-pool
provisioning, InferLine sized each stage of multi-model prediction
\emph{pipelines} under an end-to-end tail-latency SLO, pairing a
low-frequency planner with a high-frequency reactive
tuner~\cite{Crankshaw2020InferLine} --- but its stages are pre-LLM models
whose per-stage demand is the one request stream, so the question of
\emph{which signal} sizes each stage never arises. At the infrastructure level, Google's
Autopilot established history-driven (including time-series--based)
autoscaling as production practice for general
workloads~\cite{Rzadca2020Autopilot}. The common thread: the watched
signal is the service's aggregate request stream or utilization, which for
single-stage, text-only services is a reasonable proxy for its one pool's
demand.

The LLM era stressed this loop from two sides. On the supply side, cold
starts grew: loading tens of gigabytes of weights makes replica
provisioning a tens-of-seconds-to-minutes operation, and systems like
ServerlessLLM attack exactly that latency with multi-tier checkpoint
loading and locality-aware startup~\cite{Fu2024ServerlessLLM}, while
BlitzScale multicasts weights over the GPU network to make scale-up
near-live~\cite{Zhang2025BlitzScale} --- shrinking, but not eliminating,
the lead time a scaling decision must cover. On the
demand side, load became burstier and less stationary, prompting
scheduling-level responses (Llumnix reschedules requests across replicas
to absorb dynamics~\cite{Llumnix2024}), SLO-differentiated autoscalers
(Chiron scales interactive and batch LLM pools hierarchically under
separate SLOs~\cite{Patke2025Chiron}), token-granularity GPU pooling
across models in production at Alibaba~\cite{Xiang2025Aegaeon}, and,
recently, \emph{forecast-aware}
capacity management: SageServe combines a traffic-forecast model with an
ILP over multi-datacenter GPU fleets, reporting ${\sim}25\%$ GPU-hour
savings over reactive scaling~\cite{Jaiswal2025SageServe}; WarmServe
prewarms model instances on shared GPUs from five-minute peak-load
predictions, claiming order-of-magnitude TTFT improvements over reactive
autoscaling~\cite{WarmServe2025}; predictive engine-state scheduling
appears at finer granularity as well~\cite{Zhang2025SynergySched}. Dynamo's
planner compresses this whole trajectory into one artifact: introduced in
mid-2025 as an SLA planner regulating prefill and decode pool sizes on
observed TTFT/TPOT metrics, within a year it grew a second, slower
\emph{forecast-driven} loop that predicts the aggregate request rate and
mean sequence lengths with classical models (ARIMA, Kalman, Prophet) to
set replica floors under the reactive loop~\cite{Dynamo2025}. Its
reactive SLA loop is the baseline of this thesis's evaluation
\emph{[Placeholder: cross-reference the baselines chapter when it
lands]}; its predictive loop still watches a single averaged text-side
stream --- a limitation its own design notes state. Autoscalers purpose-built for
\emph{disaggregated} LLM serving are the newest arrivals, and the closest
neighbors of this thesis: ByteDance reports coordinated scaling of
heterogeneous prefill and decode pools in
production~\cite{Li2025TamingChaos}, and TokenScale sizes the two pools on
per-stage token-velocity signals rather than the request
rate~\cite{Lai2025TokenScale} --- per-stage signals, for a two-stage
\emph{text} pipeline whose stages are fed by the same token stream.

Two gaps remain. First, every system above watches (or forecasts) signals
derived from the request/token stream of a text-only fleet --- even the
per-stage token velocities of the disaggregated autoscalers are that one
stream, split by phase. Chapter~\ref{cap:trace-eda} shows that in
multimodal serving this stream
carries almost no information about the vision pool's demand
($R^2=\num{0.123}$), so grafting these autoscalers onto a
modality-disaggregated deployment reproduces exactly the missizing this
thesis quantifies. Second, their predictors are bespoke statistical
models fitted per service; none evaluates the pretrained time-series
foundation models of Section~\ref{sec:rel-forecast}, which remove the
per-signal modeling burden precisely where a multi-pool system multiplies
it.

\section{Workload characterization and traces}
\label{sec:rel-workload}

The statistical phenomena underlying this thesis's workload analysis have
deep precedent. Self-similarity and long-range dependence --- burstiness
that no aggregation window removes --- were established for Ethernet
traffic by Leland et al.~\cite{Leland1994}, for wide-area TCP traffic by
Paxson and Floyd --- whose title, \emph{the failure of Poisson modeling},
names the phenomenon~\cite{Paxson1995} --- and for Web traffic by
Crovella and Bestavros~\cite{Crovella1997}, together with the
heavy-tailed object sizes that produce them. Chapter~\ref{cap:trace-eda}
finds the same signatures (Fano factor growing by nearly three orders of
magnitude with window size, Hurst exponents up to ${\approx}0.93$) in a multimodal
inference workload's demand signals --- strongest precisely in the vision
signal.

For LLM serving specifically, public traces anchor the field's empirical
work: the Azure LLM inference traces released with
Splitwise~\cite{Patel2024Splitwise} and refreshed with
DynamoLLM~\cite{Stojkovic2025DynamoLLM}, BurstGPT's year-long
trace with its analysis of arrival burstiness~\cite{Wang2024BurstGPT},
ServeGen's characterization and generative modeling of production
workloads at Alibaba, which explicitly notes that realistic workloads
must be composed per category rather than as one
stream~\cite{Xiang2025ServeGen}, and Aegaeon's multi-month view of
bursty, long-tailed multi-model traffic in Alibaba's Model
Studio~\cite{Xiang2025Aegaeon}. All are text-centric: their request
descriptors are token counts and timestamps. The
\textit{AzureLMMInferenceDataset2025} released with
ModServe~\cite{Qiu2025ModServe} is, to our knowledge, the first public
production trace of a \emph{multimodal} inference cluster --- and
Chapter~\ref{cap:trace-eda} contributes what the autoscaling question
needs from it: the per-pool demand-signal view (decoupling, burst
anatomy, persistence gradient) that no prior characterization, including
ModServe's own, provides.

\section{Time-series forecasting for provisioning}
\label{sec:rel-forecast}

Classical operational forecasting --- exponential smoothing,
Holt--Winters, ARIMA --- is codified in Hyndman and
Athanasopoulos~\cite{Hyndman2021fpp3} and remains the default in
production autoscalers (Autopilot's recommenders, SageServe's traffic
models, WarmServe's window predictor). A recent alternative is the
pretrained \emph{time-series foundation model} (TSFM), applied zero-shot:
TimesFM, a decoder-only model pretrained on billions of
points~\cite{Das2024TimesFM}; Chronos, which tokenizes series for a
language-model backbone~\cite{Ansari2024Chronos}, and its successor
Chronos-2 with native multivariate and covariate
support~\cite{Ansari2025Chronos2}; Moirai's unified any-frequency
training~\cite{Woo2024Moirai}; Lag-Llama's probabilistic
head~\cite{Rasul2023LagLlama}; Time-MoE's billion-parameter
sparse scaling~\cite{Shi2025TimeMoE}; and Datadog's Toto family, trained
at scale on observability telemetry --- the domain closest to serving
demand signals --- released together with the BOOM benchmark of
observability series~\cite{Cohen2024Toto,Cohen2025BOOM} and scaled up in
Toto-2~\cite{Khwaja2026Toto2}. Public benchmarks such as GIFT-Eval
compare these models across broad corpora~\cite{Aksu2024GIFTEval}, but
their benchmark domains (energy, retail, weather, cloud telemetry) do not
include inference-serving demand, and the deployment-relevant questions
--- accuracy at a specific provisioning cadence, degradation across the
lead-time horizon, the value of cross-signal covariates --- are not posed
in autoscaling terms. Chapter~\ref{cap:predictors} fills exactly that
slot: a rolling-origin, provisioning-framed benchmark of three TSFMs
against classical baselines on all four demand signals of a production
multimodal trace, zero-shot, with the forecast consumed as a scaler
would consume it.

\section{Simulators and performance models}
\label{sec:rel-sim}

Systems questions of this kind are standardly answered in simulation.
Vidur simulates LLM serving at cluster scale from profiled and learned
kernel timings, validated against real deployments, and is explicitly
motivated by the cost of sweeping deployment configurations on
hardware~\cite{Agrawal2024Vidur}; LLMServingSim co-simulates the
accelerator and the serving stack, and its 2.0 release explicitly targets
heterogeneous and disaggregated serving
infrastructure~\cite{Cho2024LLMServingSim,Cho2026LLMServingSim2}; Dynamo
ships a discrete-event simulator
of its own worker pools (the \emph{mocker}), which this thesis drives
through \texttt{dynamo.replay}~\cite{Dynamo2025}; and aiconfigurator
searches deployment configurations offline by pricing operator graphs
against per-kernel performance databases collected on target
GPUs~\cite{AIConfigurator2025}. All of these model \emph{text} token
processing: none prices a vision-encoder stage, and aiconfigurator's
encoder operator had no performance data for any consumer- or
workstation-class GPU before this work
(Chapter~\ref{cap:aiconfigurator}). This thesis therefore composes its
platform from parts --- the encoder pool costed on silicon collected and
physically validated for this purpose, the language pool on Dynamo's
simulator under calibrated H100 timing --- and quantifies the fidelity of
each seam (Chapter~\ref{cap:twopool}) rather than assuming it.

\section{Positioning}
\label{sec:rel-positioning}

What separates this thesis from the landscape above is not an ingredient
but a finding. The demand signal most dangerous to reactive scaling ---
the vision pool's image rate: the burstiest of the four demand signals,
nearly invisible in the request stream a conventional autoscaler watches,
and ramping to peak within replica cold-start times in the fast half of
its bursts --- is also the most persistent, and therefore the most
forecastable, at exactly the horizons provisioning requires
(Chapters~\ref{cap:trace-eda} and~\ref{cap:predictors}). That inversion
is what makes per-pool predictive provisioning both necessary and
feasible in modality-disaggregated serving, and it is a property no prior
work, to our knowledge, states or tests.

None of the ingredients is claimed as novel here: modality
disaggregation, forecast-aware scaling, pretrained time-series
forecasters, and discrete-event serving simulation each exist, several in
production. The contribution is the demonstration that in multimodal
serving these must meet --- that no single signal can size both pools,
and that the missing signal rewards forecasting --- together with the
calibrated platform that makes the demonstration quantitative. Stated as
precedence claims, the thesis makes three, each deliberately narrow: to
our knowledge, the first per-pool demand-signal characterization of a
production multimodal inference trace; the first provisioning-framed,
zero-shot benchmark of time-series foundation models on serving demand
signals; and the first silicon-calibrated encoder cost model for a
consumer-class GPU in a deployment planner, contributed upstream.
Table~\ref{tab:rel-positioning} summarizes the landscape against the four
ingredients the claim requires. \emph{[Placeholder: extend the table's
last column once the baseline and closed-loop chapters land, citing their
headline SLA $\times$ GPU-hours result.]}

\begin{table}[htb]
\centering
\caption{Positioning against the closest related work. \checkmark\ =
addressed; $\sim$ = partially (see text); --- = not addressed.}
\label{tab:rel-positioning}
\footnotesize
\setlength{\tabcolsep}{4pt}
\begin{tabular}{lcccc}
\toprule
 & Modality & Per-pool & Forecast-driven & Calibrated \\
Work & disagg. & signals & scaling & simulation \\
\midrule
ModServe~\cite{Qiu2025ModServe} & \checkmark & $\sim$ & --- & --- \\
EPD / HydraInfer~\cite{Singh2025EPD,Dong2025HydraInfer} & \checkmark & --- & --- & --- \\
ElasticMM~\cite{Liu2025ElasticMM} & \checkmark & $\sim$ & --- & --- \\
SageServe / WarmServe~\cite{Jaiswal2025SageServe,WarmServe2025} & --- & --- & \checkmark & --- \\
TokenScale / HeteroScale~\cite{Lai2025TokenScale,Li2025TamingChaos} & --- & $\sim$ & --- & --- \\
Dynamo (planner)~\cite{Dynamo2025} & $\sim$ & --- & $\sim$ & $\sim$ \\
Vidur~\cite{Agrawal2024Vidur} & --- & --- & --- & $\sim$ \\
\textbf{This thesis} & \checkmark & \checkmark & \checkmark & \checkmark \\
\bottomrule
\end{tabular}
\end{table}
